\chapter{Classical results}
In this chapter, we present some classical results about the connections between the spectrum and the structural properties of a regular graph.
Most information is provided by the second largest and the smallest eigenvalue.

Given a $d$-regular graph $G$ on $n$ vertices, we denote its eigenvalues by 
\begin{equation*}
    d = \lambda_1 \ge \lambda_2 \ge \dots \ge \lambda_n.
\end{equation*}
Furthermore, throughout this chapter $\lambda = \lambda(G) := \max \{ |\lambda_2|, |\lambda_n | \}$ denotes the largest absolute value of an eigenvalue other than $\lambda_1 = d$.
We know that G is connected if and only if $d > \lambda_2$, and bipartite if and only if $\lambda_n = -d$. 
The difference between the largest and second largest eigenvalue of $G$, namely $d - \lambda_2$, is called the \textit{spectral gap} \mbox{of $G$}.
We will see that if the spectral gap is large, that is, $\lambda_2$ is small compared to $d$, $G$ has good expansion properties. 
A small value of $-\lambda_n$ implies that the size of an independent set in $G$ cannot be too large.
In the case where both $\lambda_2$ and $-\lambda_n$ are small, the edge distribution of the graph even resembles that of a random graph.  

\section{Graph expansion and Cheeger's inequality}
    \label{section:cheeger}
We now introduce the concept of graph expansion, following the notation and terminology used in the survey \cite{expandergraphs}. 
For $S, T \subseteq V(G)$, denote the set of edges from $S$ to $T$ by
\begin{equation*}
    E(S, T) := \{ (x, y) \in S \times T : xy \in E(G) \}
\end{equation*}
and let $e(S, T) := |E(S, T)|$. 
Note that $e(S, T)$ counts directed edges; in particular, edges within $S \cap T$ are counted twice.
Furthermore, define $E(S)$ as the set of all (undirected) edges with both endvertices in $S$, and $e(S):= |E(S)|$.

\begin{definition}
    Let $G$ be a graph on $n$ vertices, and let $S \subseteq V(G)$.
    The \textit{edge boundary} of $S$, denoted $\partial S$, is the set
    \begin{equation*}
        \partial S := E(S, \overline{S})
    \end{equation*}
    of all edges from a vertex in $S$ to a vertex in $\overline{S} := V(G) \setminus S$.
    The \textit{(edge) expansion ratio} $h(G)$ of $G$ is defined as 
    \begin{equation*}
        h(G) := \min \bigg \{ \frac{|\partial S|}{|S|} : S \subseteq V(G),\, 0 < |S|\; \le \frac{n}{2} \bigg \}.
    \end{equation*}
\end{definition}

Note that the expansion ratio of a graph $G$ is positive if and only if $G$ is connected. 

\begin{definition}
    Let $\varepsilon > 0$. A family of $d$-regular graphs $(G_i)_{i \in \N}$ of size increasing with $i$ is called a \textit{family of} $\varepsilon$-\textit{expander graphs} if
    \begin{equation*}
        h(G_i) \ge \varepsilon, \ \ i \in \N.
    \end{equation*}
\end{definition}

Qualitatively speaking, a family of expander graphs consists of sparse graphs with the property that, for any set of vertices $S$ that is "not too large", the number of edges leaving $S$ is bounded from below by $\varepsilon |S|$ for some constant $\varepsilon > 0$. 
This implies that a good expander cannot be split into two large parts by removing only a few edges, a desirable property, for example, in the design of robust communication networks. 
Since their introduction in the 1970s, expander graphs have found various applications in mathematics and computer science, e.g. in the construction of good error-correcting codes, the theory of metric embeddings (see \cite{expandergraphs}) and sorting networks (see \cite{aks}).

As the properties of sparseness and good connectivity are seemingly conflicting, it may not be immediately clear that infinite families of $d$-regular expander graphs even exist.
However, in the last decades, numerous constructions – both explicit and probabilistic – have been developed.
This shall not be the focus of this thesis; for a more thorough introduction to expander graphs, we refer the reader to \cite{expandergraphs}.

We should note that there are other interesting notions of graph expansion as well.
The most notable variation of the above definition is vertex expansion, where instead of the edge boundary, the neighbourhood of a vertex set is considered.
We will briefly encounter this notion of expansion in Chapter 4.

Cheeger's inequality, one of the most fundamental results of spectral graph theory, relates the edge expansion ratio of a graph to its spectral gap. 

\begin{theorem}[Cheeger's inequality] \label{thm:cheeger}
    Let $G$ be a $d$-regular graph on $n$ vertices with second-largest eigenvalue $\lambda_2$. Then 
    \begin{equation*}
        \frac{d - \lambda_2}{2} \le h(G) \le \sqrt{2d (d - \lambda_2)}.
    \end{equation*}
\end{theorem}

The inequality tells us that a large spectral gap implies a high expansion ratio, and vice versa. 
The first versions of this theorem have been proved in the 1980's (see \cite{dodziuk, aloneigenvaluesandexpanders}), inspired by an analogous theorem of Cheeger \cite{cheeger} in the setting of Riemannian manifolds.
Since then, the result has been refined and generalised several times.
For instance, a similar result exists for the \textit{conductance}, an analogue of the expansion ratio for weighted graphs (see e.g. \cite[Chapter 21]{spielman}).

Recall the definition of the Laplace matrix $L = D - A$ of $G$, where $D$ is the degree diagonal matrix.
In the literature, Cheeger's inequality is often stated in terms of the second smallest eigenvalue \mbox{of $L$} (the smallest one is always zero).
This way, it can naturally be generalised to the case of irregular graphs.
In the case of a $d$-regular graph, the definition of the Laplace matrix simplifies to $L = dI - A$, and it is clear that the second-smallest eigenvalue of $L$ is precisely the spectral gap of $G$, i.e. \mbox{$d - \lambda_2$}.

\begin{remark}
    Both the lower and upper bounds are essentially tight:
    \begin{enumerate}
        \item The lower bound is met by the $d$-dimensional hypercube $Q_d$. 
            By Example \ref{ex:special_spectra}, the spectral gap is $d - \lambda_2 = 2$. 
            On the other hand, we have $h(Q_d) \le 1$, since $|\partial S| / |S|\; = 1$ for any $(d-1)$-dimensional subcube $S$ of $Q_d$. 

        \item The cycle $C_n$ shows that the upper bound is tight up to a constant factor.
            Here 
            \begin{equation*}
                h(C_n) = \frac{2}{\lfloor n / 2 \rfloor} = \Theta \Big (\frac{1}{n} \Big),
            \end{equation*}
            attained by the half cycle on $\lfloor n/2 \rfloor$ vertices.
            Recalling again our discussion in Example \ref{ex:special_spectra}, we have $\lambda_2 = 2 \cos ( {2 \pi}/{n})$ (and, of course, $\lambda_1 = 2$).
            Now, by the power series expansion of the cosine, 
            \begin{equation*}
                \cos \Big( \frac{2 \pi}{n} \Big) = 1 - \frac{2 \pi^2}{n^2} + O \Big( \frac{1}{n^4} \Big) = 1 - \Theta \Big( \frac{1}{n^2} \Big),
            \end{equation*}
            and hence the spectral gap is $\Theta ({1}/{n^2})$.
    \end{enumerate}
\end{remark}

For the proof of Cheeger's inequality, we will follow \cite{expandergraphs} and show the lower and upper bounds of the theorem separately.

\begin{proof}[Proof of the lower bound]
    Recall that the principal eigenvector of $G$ is $\bm{1}$ because $G$ is regular.
    By Rayleigh's principle, 
    \begin{equation*}
        \lambda_2 \ge \frac{\bm{u}^T A \bm{u}}{\bm{u}^T \bm{u}}
    \end{equation*}
    for any vector $\bm{u} \perp \bm{1}$. 
    We are going to exhibit a vector $\bm{u} \perp \bm{1}$ with a Rayleigh quotient of at least $d - 2 h(G)$, hence proving $\lambda_2 \ge d - 2 h(G)$.
    Let $S \subseteq V(G)$ be a subset of the vertices satisfying $h(G) = |\partial S|/|S|$ and $|S|\, \le n/2$ and consider $\bm{u} = |\overline{S}| \cdot \bm{1}_S - |S| \cdot \bm{1}_{\overline{S}}$.
    Clearly, $\bm{u}^T \bm{1} = 0$. Furthermore,
    \begin{equation*}
        \bm{u}^T \bm{u} = 
        %|\overline{S}|^2 \bm{1}_S^T \bm{1}_S - |\overline{S}| |S| \bm{1}_S^T \bm{1}_{\overline{S}} - |S| |\overline{S}| \bm{1}_{\overline{S}}^T \bm{1}_S + |S|^2 {\bm{1}_{\overline{S}}}^T \bm{1}_{\overline{S}}
            |\overline{S}|^2 |S|\, -\, |\overline{S}| |S| \cdot 0 - |S| |\overline{S}| \cdot 0 + |S|^2 |\overline{S}|\;
        = (|\overline{S}| \,+\, |S|) |S| |\overline{S}|  = n |S| |\overline{S}| 
    \end{equation*}
    and
    \begin{equation*}
        \bm{u}^T A \bm{u} = 2|\overline{S}|^2 e(S)\,+\, 2|S|^2 e(\overline{S})\, -\, 2 |\overline{S}| |S| e(S, \overline{S}).
    \end{equation*}
    Since $G$ is $d$-regular, we have $2 e(S) = d|S|\, -\, e(S, \overline{S})$ and $2 e(\overline{S}) = d |\overline{S}|\, -\, e(S, \overline{S})$. It follows that 
    \begin{align*}
        \bm{u}^T A \bm{u} &= d|\overline{S}|^2 |S| \,-\, |\overline{S}|^2 e(S, \overline{S}) + d|S|^2 |\overline{S}| \,-\, |S|^2 e(S, \overline{S}) - 2|\overline{S}| |S| e(S, \overline{S}) \\
        &= d|S||\overline{S}| (|S| \,+\, |\overline{S}|) - (|S|^2 \,+\, 2 |\overline{S}| |S| \,+\, |\overline{S}|^2) e(S, \overline{S}) \\
        &= nd |S| |\overline{S}|\,-\, n^2 e(S, \overline{S}).
    \end{align*}
    Using that $h(G) = |e(S, \overline{S})|/|S|$ and $|\overline S|\; \ge n/2$, we conclude
    \begin{equation*}
        \lambda_2 \ge \frac{\bm{u}^T A \bm{u}}{\bm{u}^T \bm{u}} = \frac{nd |S| |\overline{S}| \,-\, n^2 e(S, \overline{S})}{n |S| |\overline S|} = d - \frac{n e(S, \overline S)}{|S| |\overline S|} \ge d - 2h(G).
    \end{equation*}
\end{proof}

\begin{proof}[Proof of the upper bound]
    This part is known as the difficult direction. 
    Let $\bm{x}$ be a $\lambda_2$-eigenvector of $G$. 
    We want to partition $V(G)$ according to the sign of $\bm{x}$ (which has positive and negative entries since $\bm{x} \perp \bm{1}$).
    The intuition here is that this partition somehow constitutes a sparse cut in $G$, that is, a partition of the vertex set in two large parts such that the number of edges joining both parts is small.
    Define $\bm{y} = \bm{x}^{+}$ (that is, $\bm{y}_v = \max \{ \bm{x}_v, 0 \}$ for $v \in V(G)$) and 
    \begin{equation*}
        V^{+} = \supp(\bm{y}) = \{ v \in V(G) : \bm{y}_v > 0 \}.
    \end{equation*}
    Without loss of generality, we can assume that $|V^+|\; \le n/2$ (otherwise consider $-\bm{x}$, which is a $\lambda_2$-eigenvector as well).

    Let $L$ be the Laplace matrix of $G$. Since $G$ is $d$-regular, we have $L = dI - A$.
    We wish to prove our inequality by establishing lower and upper bounds for the Rayleigh quotient $\bm{y}^T L \bm{y}/\bm{y}^T \bm{y}$.

    \begin{claim}
        \begin{equation*}
            \frac{\bm{y}^T L \bm{y}}{\bm{y}^T \bm{y}} \le d - \lambda_2.
        \end{equation*}
    \end{claim}
    
    \begin{proof}
        For $v \in V^+$, we can write
        \begin{align*}
            (L\bm{y})_v &= d \bm{y}_v - \sum \limits_{w \in V(G)} A_{vw} \bm{y}_w \\
                &= d \bm{x}_v - \sum \limits_{w \in V^+} A_{vw} \bm{x}_w \\
                &\le d \bm{x}_v - \sum \limits_{w \in V(G)} A_{vw} \bm{x}_w 
                = (d - \lambda_2) \cdot \bm{x}_v.
        \end{align*}
        Since $\bm{y}_v = 0$ for $v \notin V^+$, it follows that 
        \begin{align*}
            \bm{y}^T L \bm{y} &= \sum \limits_{v \in V(G)} \bm{y}_v (L\bm{y})_v \\
                &\le (d - \lambda_2) \sum \limits_{v \in V^+} \bm{x}_v^2 \\ 
                &= (d - \lambda_2) \sum \limits_{v \in V(G)} \bm{y}_v^2 = (d - \lambda_2) \cdot \bm{y}^T \bm{y},
        \end{align*}
        establishing the claim.
    \end{proof}

    \begin{claim} \label{claim:cheeger-lower}
        \begin{equation*}
            \frac{h(G)^2}{2d} \le \frac{\bm{y}^T L \bm{y}}{\bm{y}^T \bm{y}}.
        \end{equation*}
    \end{claim}

    \begin{proof}
        For ease of notation, let $h = h(G)$.
        Define
        \begin{equation*}
            B = \sum \limits_{v \sim w} | \bm{y}_v^2 - \bm{y}_w^2 |,
        \end{equation*}
        where the summation is meant to be over each edge $vw \in E(G)$.
        We are going to show that 
        \begin{equation*}
            h \cdot \|\bm{y} \|^2\, \le B \le \sqrt{2d} \cdot \|\bm{y} \| \cdot (\bm{y}^T L \bm{y})^{\frac{1}{2}},
        \end{equation*}
        which yields the desired inequality by squaring.

        We start by proving the upper bound on $B$. 
        By the Cauchy--Schwarz inequality, we have
        \begin{align*}
            B &= \sum \limits_{v \sim w} | \bm{y}_v^2 - \bm{y}_w^2 | \\
                &= \sum \limits_{v \sim w} |\bm{y}_v + \bm{y}_w| | \bm{y}_v - \bm{y}_w|  \\ 
                &\le \Bigg( \sum \limits_{v \sim w} (\bm{y}_v + \bm{y}_w)^2 \Bigg)^{\frac{1}{2}}  \Bigg( \sum \limits_{v \sim w} (\bm{y}_v - \bm{y}_w)^2 \Bigg)^{\frac{1}{2}} .
        \end{align*}
        By Proposition \ref{prop:laplacian-quadratic-form}, the second factor is
        \begin{align*}
            \Bigg ( \sum \limits_{v \sim w} (\bm{y}_v - \bm{y}_w)^2 \Bigg)^{\frac{1}{2}} = (\bm{y}^T L \bm{y})^{\frac{1}{2}}.
        \end{align*}
        Evaluating the first factor gives
        \begin{align*}
            \Bigg( \sum \limits_{v \sim w} (\bm{y}_v + \bm{y}_w)^2 \Bigg)^{\frac{1}{2}} &\le \Bigg( \sum \limits_{v \sim w} 2 (\bm{y}_v^2 + \bm{y}_w^2) \Bigg)^{\frac{1}{2}} \\
                &= \Bigg( 2d \cdot \sum \limits_{v \in V(G)} \bm{y}_v^2 \Bigg)^{\frac{1}{2}} = \sqrt{2d} \cdot \| \bm{y} \|.
        \end{align*}
        The first inequality follows from the fact that  
        \begin{equation*}
            0 \le (\bm{y}_v - \bm{y}_w)^2 = \bm{y}_v^2 - 2 \bm{y}_v \bm{y}_w + \bm{y}_w^2,
        \end{equation*}
        so $\bm{y}_v^2 + \bm{y}_w^2 \ge 2 \bm{y}_v \bm{y}_w$.
        For the middle equality we have used that each $\bm{y}_v^2$ occurs precisely $d$ times in the sum, namely for every edge incident with $v$.
        Combining these inequalities, we conclude 
        \begin{equation*}
             B \le \sqrt{2d} \cdot \|\bm{y} \|  \cdot (\bm{y}^T L \bm{y})^{\frac{1}{2}}.
        \end{equation*}

        It remains to show the lower bound on $B$. 
        For the rest of the proof, we assume that the vertices in $V(G)$ are labelled by $1, \dots, n$ such that $\bm{y}_1 \ge \dots \ge \bm{y}_n$.
        % In particular, $\bm{y}_{i} = 0$ for all $i > |V^+|$.
        We have
        \begin{align*}
            B &= \sum \limits_{\substack{v \sim w \\ v < w}} (\bm{y}_v^2 - \bm{y}_w^2) \\
                    &= \sum \limits_{\substack{v \sim w \\ v < w}} \sum \limits_{i = v}^{w - 1} (\bm{y}_i^2 - \bm{y}_{i + 1}^2) 
                = \sum \limits_{i = 1}^{n - 1} (\bm{y}_i^2 - \bm{y}_{i + 1}^2) \cdot e([i], \overline{[i]}).% \\ 
        \end{align*}
        For the last equality, observe that the summand $(\bm{y}_i^2 - \bm{y}_{i + 1}^2)$ appears for every edge $vw$ with $v \le i < w$, so precisely $e([i], \overline{[i]})$ times in total. 
        As $\bm{y}_{i} = 0$ for all $i > |V^+|$, it follows that
        \begin{align*}
            B &= \sum \limits_{i \in V^+} (\bm{y}_i^2 - \bm{y}_{i + 1}^2) \cdot e([i], \overline{[i]}) \\ 
                &\ge h \cdot \sum \limits_{i \in V^+} (\bm{y}_i^2 - \bm{y}_{i + 1}^2) \cdot i 
                = h \cdot \sum_{i \in V^+} \bm{y}_i^2 = h \|\bm{y}\|^2.
        \end{align*}
        The inequality holds because $h \le e([i], \overline{[i]}) / i$ for $i \in V^+$. 
        (Here the assumption that $|V^+|\; \le n/2$ comes to fruition.)
        Altogether, this establishes Claim \ref{claim:cheeger-lower}.
    \end{proof}

    Combining the claims, we obtain
    \begin{equation*}
        \frac{h(G)^2}{2d} \le \frac{\bm{y}^T L \bm{y}}{\bm{y}^T \bm{y}} \le d - \lambda_2,
    \end{equation*}
    so $h(G) \le \sqrt{2d (d - \lambda_2)}$. This concludes the proof of Cheeger's inequality.
\end{proof}

Cheeger's inequality implies that expansion and spectral gap are qualitatively equivalent, so instead of constructing graphs with a high expansion ratio directly, one can also search for graphs with a large spectral gap.
It is therefore of interest how large this gap can become.
Clearly, the answer depends on whether $d$ grows with $n$ as $n \rightarrow \infty$. 
The extremal example of the complete graph $K_n$, where $d = n - 1$, demonstrates that in this case one cannot expect a very meaningful result:
$K_n$ has spectrum $n-1,  -1^{(n - 1)}$, so the spectral gap is $n$.

However, if $d$ is fixed and $n$ grows large, the Alon--Boppana bound provides a satisfactory answer.
\begin{theorem}[Alon--Boppana]
    For every $d$-regular graph $G$ on $n$ vertices, 
    \begin{equation*}
        \lambda_2(G) \ge 2\sqrt{d - 1} - o(1),
    \end{equation*}
    where the term $o(1)$ tends to $0$ for fixed $d$ as $n \rightarrow \infty$.
\end{theorem}

Two different proofs for this result are given in \cite[Section 5]{expandergraphs}. 
The following weaker statement is much more easily obtained, and it demonstrates that $\sqrt{d}$ is the right order of magnitude for an asymptotic lower bound on $\lambda = \max \{ |\lambda_2|, |\lambda_n| \}$.

\begin{proposition}
    For every $d$-regular graph $G$ on $n$ vertices, 
    \begin{equation*}
        \lambda \ge \sqrt{d} \cdot (1 - o(1)),
    \end{equation*}
    where the term $o(1)$ tends to $0$ for fixed $d$ as $n \rightarrow \infty$.
\end{proposition}

\begin{proof}
    Recall that $\tr A^2 = 2 e(G)$. 
    Since $G$ is $d$-regular, this number equals $nd$, and we have
    \begin{equation*}
        nd = \tr A^2 = \sum \limits_{i = 1}^n \lambda_i^2 \le d^2 + (n - 1)\lambda^2.
    \end{equation*}
    It follows that
    \begin{equation*}
        \lambda^2 \ge d \cdot \frac{n - d}{n - 1},
    \end{equation*}
    so $\lambda \ge \sqrt{d} \cdot (1 - o(1))$, as desired.
\end{proof}

We end this section with a brief mention of graphs that almost attain the Alon--Boppana bound.
\begin{definition}
    A $d$-regular graph $G$ is called a \textit{Ramanujan graph} if each eigenvalue $\lambda \ne \pm d$ satisfies $|\lambda|\; \le 2\sqrt{d - 1}$.
\end{definition}

There are many examples of small Ramanujan graphs, such as the complete \sloppy \mbox{graph $K_{d +1}$}.
It is, however, very challenging to construct (for a fixed $d$) infinite families of $d$-regular graphs that are all Ramanujan.
Indeed, in light of the Alon--Boppana bound, the condition for a graph to be Ramanujan is strongest possible, since for every $\varepsilon > 0$, any sufficiently large $d$-regular graph satisfies $\lambda_2 > 2\sqrt{d - 1} - \varepsilon$.

A distinction is made between non-bipartite Ramanujan graphs and bipartite Ramanujan graphs (i.e. those having $-d$ as an eigenvalue).
It is known that arbitrarily large $(q^k + 1)$-regular non-bipartite Ramanujan graphs exist for any prime number $q$ and $k \ge 1$ (see \cite{morgensternramanujangraphs}).
Furthermore, Marcus, Spielman, and Srivastava \cite{interlacingfamilies} gave a probabilistic construction of infinite families of $d$-regular bipartite Ramanujan graphs for any integer $d \ge 3$.
In particular, these results show that the Alon--Boppana bound is tight.
Whether infinite families of $d$-regular non-bipartite Ramanujan graphs exist for any integer $d \ge 3$ remains an open problem.

\section{Hoffman's bound} \label{section:hoffman-bound}
The next result due to Hoffman provides a bound on the maximum size of an independent set in terms of the smallest eigenvalue $\lambda_n$.
\begin{theorem}[Hoffman's bound]
    Let $G$ be a $d$-regular graph. Then 
    \begin{equation*}
        \alpha(G) \le n \frac{-\lambda_n}{d - \lambda_n}.
    \end{equation*}
    Moreover, if an independent set $U$ attains this bound, then each vertex in $V(G) \setminus U$ has precisely $-\lambda_n$ neighbours in $U$. 
\end{theorem}

Note that in a non-empty graph, $\lambda_n < 0$.

We follow the proof given in \cite[Theorem 3.5.2]{brouwerhaemers}, which uses interlacing on a particular quotient matrix.
\begin{proof}
    Let $U \subseteq V(G)$ be an independent set of size $\alpha := \alpha(G)$.
    Consider a partition of $V(G)$ according to $U$ and $V(G) \setminus U$.
    The corresponding quotient matrix has the form
    \begin{equation*}
        B = \begin{bmatrix}
            0 & d \\ 
            \frac{\alpha d}{n - \alpha} & d - \frac{\alpha d}{n - \alpha}
        \end{bmatrix}.
    \end{equation*}
    By multiplying $B$ with $\bm{1}$, we see that $\bm{1}$ is an eigenvector of $B$ for the eigenvalue $\mu_1 = d$.
    Since $\tr B = d - \frac{\alpha d}{n - \alpha}$, the second eigenvalue must be $\mu_2 = -\frac{\alpha d}{n - \alpha}$.
    By interlacing (Proposition \ref{prop:partition-interlacing}) it follows that 
    \begin{equation*}
        -\frac{\alpha d}{n - \alpha} = \mu_2 \ge \lambda_n.
    \end{equation*}
    Solving for $\alpha$, we obtain the desired inequality:
    \begin{align*}
        -\frac{\alpha d}{n - \alpha} \ge \lambda_n \ &\Longleftrightarrow \ \alpha d \le -\lambda_n (n - \alpha)\\ 
            &\Longleftrightarrow \ \alpha (d - \lambda_n) \le -n \lambda_n \\ 
            &\Longleftrightarrow \ \alpha \le n \frac{- \lambda_n}{d - \lambda_n}.
    \end{align*}

    Now assume that equality holds for $U$, so we must have $\mu_2 = \lambda_n$.
    Since also $d = \lambda_1 = \mu_1$, the interlacing is tight, and hence the partition is equitable. 
    This means that the term $\frac{\alpha d}{n - \alpha} = -\lambda_n$ in the quotient matrix is precisely the number of neighbours that each vertex of $V(G) \setminus U$ has in $U$.
\end{proof}

Hoffman's bound can be used to prove the well-known Erd\H os--Ko--Rado theorem.

\begin{theorem}[Erd\H os--Ko--Rado, \cite{erdoskorado}]
    \label{thm:erdoskorado}
    Let $v, r$ be positive integers with $v \ge 2r$.
    Suppose $\mathcal A$ is a family of distinct $r$-element subsets of an $v$-element set $V$, in which any two subsets have at least one common element.
    Then
    \begin{equation*}
        |\mathcal A |\;\le \binom{v - 1}{r - 1}.
    \end{equation*}
\end{theorem}

Let $v \ge 2r$ as in the theorem, and let $V$ be a set of size $v$.
The \textit{Kneser graph} $K(v, r)$ is the graph with vertex set $\binom{V}{r}$, i.e. the set of $r$-element subsets of $V$, in which two vertices are adjacent if and only if the corresponding subsets are disjoint.
Clearly, we have $|K(v, r)| = \binom{v}{r}$, and each vertex has degree $\binom{v - r}{r}$.

An independent set in $K(v, r)$ corresponds to a family of $r$-subsets of $V$ such that any two subsets share at least one element. 
Hence, the Erd\H os--Ko--Rado theorem has an equivalent succinct graph-theoretic formulation: 
\begin{equation*}
    \alpha(K(v, r)) \le \binom{v - 1}{r - 1}.
\end{equation*}

Using Hoffman's bound, one can give a very short proof for this inequality.
\begin{proof}[Proof of Theorem \ref{thm:erdoskorado}]
    It is known (see e.g. \cite[Section 9.4]{godsilroyle}) that the least eigenvalue of the Kneser graph $K(v, r)$ is 
    \begin{equation*}
        - \binom{v - r - 1}{r - 1}.
    \end{equation*}
    Hence, by Hoffman's bound, 
    \begin{align*}
        \alpha(K(v, r)) &\le \binom{v}{r} \frac{\binom{v - r - 1}{r - 1}}{\binom{v - r}{r} + \binom{v - r - 1}{r - 1}} 
            = \frac{\binom{v}{r}}{1 + \frac{\binom{v - r}{r}}{\binom{v - r - 1}{r - 1}}} 
            = \frac{\binom{v }{r}}{1 + \frac{v - r}{r}}
            = \frac{r}{v} \binom{v}{r} = \binom{v - 1}{r - 1}.
    \end{align*}
\end{proof}

It is worth noting that the bound of the Erd\H os--Ko--Rado theorem is tight:
For any fixed $a \in V$, the family $\mathcal A$ of all $r$-subsets of $V$ that contain $a$ has size precisely $\binom{v - 1}{r - 1}$.

Since the problem of determining the independence number $\alpha(G)$ of a graph $G$ is NP-hard, any easily computable bound on $\alpha(G)$ is potentially also of practical interest. 
The above observation on the tightness of the Erd\H os--Ko--Rado theorem suggests that Hoffman's bound can indeed often provide a good approximation for $\alpha(G)$.

\section{Expander Mixing Lemma}
For a $d$-regular graph $G$, let $\lambda = \lambda(G)$ denote the largest absolute value of an eigenvalue other than $\lambda_1 = d$, that is, $\lambda = \max \{ |\lambda_2|, |\lambda_n| \}$.
The following well-known lemma tells us that if $\lambda$ is small, then the edges of $G$ are "evenly distributed" over the graph.

\begin{lemma}[Expander mixing lemma]
    For all $S, T \subseteq V(G)$ the following holds:
    \begin{equation*}
        \bigg| e(S, T) - \frac{d |S||T|}{n} \bigg| \le \lambda \sqrt{|S| |T|}.
    \end{equation*}
\end{lemma}

The quantity $d|S||T|/n$ is the expected number of edges between $S$ and $T$ in a randomly chosen $d$-regular graph of order $n$.
Hence, a small value of $\lambda$ implies that the actual number of edges between $S$ and $T$ does not deviate too much from the expected number, so the edge distribution is similar to that of a random graph.

\begin{remark}
    Bipartite graphs can have a large spectral gap, but they always satisfy $\lambda_n = -d$.
    Since a bipartite graph hardly looks like a random $d$-regular graph, it is plausible that the term $\lambda \sqrt{|S||T|}$ cannot be replaced by $\lambda_2 \sqrt{|S||T|}$.
    % For example, consider the complete bipartite graph $K_{\frac{n}{2}, \frac{n}{2}}$.
    % Its second-largest eigenvalue $\lambda_2$ is $0$ (see Example \ref{ex:special_spectra}).
    % However, if we choose $S$ and $T$ to be the two partition classes of $K_{\frac{n}{2}, \frac{n}{2}}$, we obtain
    % \begin{equation*}
    %     \bigg| e(S, T) - \frac{d |S||T|}{n} \bigg| = \bigg | \frac{n^2}{4} - \frac{n^2}{8} \bigg| = \frac{n^2}{8} > 0.
    % \end{equation*}
    For example, consider a \mbox{$d$-regular} bipartite Ramanujan graph $G$ with $d \ge 15$. 
    Choose $S$ and $T$ to be the two partition classes of $G$.
    Since $G$ is regular, $|S|\, = |T|\, = n/2$, and we obtain %hence $\lambda_2 \sqrt{|S| |T|} \le \sqrt{d - 1} \cdot n$. 
    % However,
    \begin{equation*}
        \bigg| e(S, T) - \frac{d |S||T|}{n} \bigg| = \bigg| \frac{dn}{2} - \frac{d n^2}{4n} \bigg| = \frac{dn}{4} > \sqrt{d - 1} \cdot n \ge \lambda_2 \sqrt{|S| |T|},
    \end{equation*}
    where the penultimate inequality holds because $d \ge 15$.
    This illustrates that a small value of $\lambda_2$ alone is not enough to guarantee this sort of "pseudorandomness".
\end{remark}

For the proof of the expander mixing lemma, we shall use the following small observation from linear algebra:
\begin{observation} \label{observation:scalarproducts}
    Let $\bm{v}_1, \dots, \bm{v}_n$ be an orthonormal basis of $\R^n$, and let $\bm{x} \in \R^n$. 
    If $\bm{x} = \alpha_1 \bm{v}_1 + \dots + \alpha_n \bm{v}_n$, then $\alpha_i = \langle \bm{x}, \bm{v}_i \rangle$.
    Furthermore, $\| \alpha \|\; = \| \bm{x} \|$.
\end{observation}

\begin{proof}
    Using the orthonormality of the $\bm{v}_i$, we have
    \begin{align*}
        \sum \limits_{i = 1}^n \alpha_i \bm{v}_i 
            &= \sum \limits_{i = 1}^n \sum \limits_{j = 1}^n \alpha_j \langle \bm{v}_j, \bm{v}_i \rangle \bm{v}_i 
            = \sum \limits_{i = 1}^n \left \langle \sum \limits_{j = 1}^n \alpha_j \bm{v}_j, \bm{v}_i \right \rangle \bm{v}_i 
                = \sum \limits_{i = 1}^n \langle \bm{x}, \bm{v}_i \rangle \bm{v}_i,
    \end{align*}
    and therefore $\alpha_i = \langle \bm{x}, \bm{v}_i \rangle$ for $i \in [n]$. 
    As for the second part, 
    \begin{equation*}
        \| \bm{x} \|^2 \;= \langle \bm{x}, \bm{x} \rangle 
            % = \left \langle \sum \limits_{i = 1}^n \alpha_i \bm{v}_i, \sum \limits_{i = 1}^n \alpha_i \bm{v}_i \right \rangle
            = \sum \limits_{i = 1}^n \sum \limits_{j = 1}^n \langle \alpha_i \bm{v}_i, \alpha_j \bm{v}_j \rangle  = \sum \limits_{i = 1}^n \alpha_i ^2 = \| \alpha \|^2,
    \end{equation*}
    where the orthonormality was used for the penultimate equality.
\end{proof}

\begin{proof}[Proof of the expander mixing lemma]
    Once again, we follow the proof given in \cite{expandergraphs}.
    Let $\bm{1}_S, \bm{1}_T$ be the characteristic vectors of $S$ and $T$ (i.e., the $v$-th entry of $\bm{1}_S$ is one if $v \in S$ and zero otherwise).
    Let $\bm{v}_1, \dots, \bm{v}_n$ be an orthonormal basis of eigenvectors (such that $A \bm{v}_i = \lambda_i \bm{v}_i$).
    Since $G$ is regular, we can assume $\bm{v}_1 = \bm{1} / \sqrt{n}$.
    Write $\bm{1}_S = \sum_i \alpha_i \bm{v}_i$ and $\bm{1}_T = \sum_j \beta_j \bm{v}_j$. Then
    \begin{align*}
        e(S, T) = \bm{1}_S ^T A \bm{1}_T 
                &= \Bigg( \sum \limits_{i = 1}^n \alpha_i \bm{v}_i \Bigg)^T A \Bigg( \sum \limits_{j = 1}^n \beta_j \bm{v}_j \Bigg) \\
                &= \Bigg( \sum \limits_{i = 1}^n \alpha_i \bm{v}_i \Bigg)^T \Bigg(\sum \limits_{i = 1}^n \beta_j \lambda_j \bm{v}_j \Bigg) 
                = \sum \limits_{i = 1}^n \lambda_i \alpha_i \beta_i,
    \end{align*}
    where for the last equality we have used that the $\bm{v}_i$ are orthonormal. 
    By Observation \ref{observation:scalarproducts}, we have \mbox{$\alpha_1 = \langle \bm{1}_S, {\bm{1}}/{\sqrt{n}} \rangle = {|S|}/{\sqrt{n}}$}, and similarly, $\beta_1 = {|T|}/{\sqrt{n}}$.
    Since $\lambda_1 = d$, it follows that
    \begin{equation*}
        e(S, T) = d \cdot \frac{|S||T|}{n} + \sum \limits_{i = 2}^n \lambda_i \alpha_i \beta_i.
    \end{equation*}
    Finally, by the definition of $\lambda$ and using Cauchy--Schwarz, we obtain
    \begin{align*}
        \bigg| e(S, T) - \frac{d |S||T|}{n} \bigg| &= \bigg|  \sum \limits_{i = 2}^n \lambda_i \alpha_i \beta_i \bigg|  
            \le \lambda \sum \limits_{i = 2}^n |\alpha_i \beta_i| \\
            &\le \lambda \| \alpha \| \| \beta \| \; = \lambda \| \bm{1}_S \| \| \bm{1}_T \|\; = \lambda \sqrt{|S| |T|},
    \end{align*}
    completing the proof.
\end{proof}

Bilu and Linial \cite{bilulinial} showed that a converse of the expander mixing lemma holds as well.
\begin{lemma}
    Let $G$ be a $d$-regular graph satisfying 
    \begin{equation*}
        \bigg| e(S, T) - \frac{d |S||T|}{n} \bigg| \le \alpha \sqrt{|S| |T|},
    \end{equation*}
    for any two disjoint sets $S, T \subseteq V(G)$ and for some $\alpha > 0$.
    Then 
    \begin{equation*}
        \lambda(G) \le O (\alpha \cdot (1 + \log (d / \alpha))).
    \end{equation*}
\end{lemma}
